% GENERATED FILE. Edit canonical lessons, then run npm run generate:books.
% canonical-source-hash: fnv1a64:6e5b5f49329d3d38
% canonical-lessons: JA-C325-shitagau, JA-C325-sakarau, JA-C325-uketomeru, JA-C325-kekkyoku, JA-C325-jiyuu

\chapter{Going Along or Pushing Back}
\label{ch:ja-going-along-or-pushing-back}
\hlchaptermodality{325}

\begin{chapteropening}
\textbf{By the end of this chapter:} \emph{I can say whether I will go along with a decision, and why.}

Complete the last lesson of chapter 325: I can say whether I will go along with a decision, and why.
\end{chapteropening}

\section[\texorpdfstring{\ja{したがう} (shitagau)}{したがう (shitagau)}]{\ja{したがう} (shitagau) — to obey; to follow}

\label{lesson:JA-C325-shitagau}



\begin{quote}
Before the new one: say the Japanese for to criticize, then the Japanese for an excuse.
\end{quote}

\subsection*{You'll want to know: \ja{したがう}}
\textbf{\ja{したがう}} — \emph{shitagau} — \textquotedblleft{}to obey; to follow\textquotedblright{}.

\begin{quote}
\textbf{\ja{せんせいの} \ja{いけんに} \ja{したがった} \ja{ほうが} \ja{いいです。}} — \emph{sensei no iken ni shitagatta hō ga ii desu} — \textquotedblleft{}You should follow the teacher's advice.\textquotedblright{}
\end{quote}

Polite \textbf{\ja{したがいます}}, \emph{shitagaimasu}; a godan verb. Past (polite): \textbf{\ja{がくせいは} \ja{したがいました。}}, \emph{gakusei wa shitagaimashita}.

\begin{cousinweb}
\textbf{\ja{したがう}} (\emph{shitagau}) is a native Japanese word (wago), inherited rather than borrowed from Chinese, and its older origin is not settled. Its kanji was attached to it later as a kun'yomi, a native reading given to a Chinese character (JMdict, KANJIDIC2).
\end{cousinweb}

\begin{grammarlens}[title={What you've built}]
One more word for giving reasons and opinions.
\end{grammarlens}

\subsection*{Your turn}
\begin{itemize}
\raggedright
  \item \emph{Say it:} \emph{shitagau}
  \item \emph{Say it:} \emph{shitagau}, once more
  \item \emph{Say it:} \emph{sensei no iken ni shitagatta hō ga ii desu} — You should follow the teacher's advice.
  \item \emph{Recall:} say the Japanese for to criticize, then the Japanese for an excuse, then say \emph{shitagau} again
\end{itemize}

\begin{tcolorbox}[breakable,colback=teal!4,colframe=teal!35!black,title={Before you move on}]
What does \ja{したがう} mean? (\textquotedblleft{}To obey; to follow\textquotedblright{}.) Say it once more.
\end{tcolorbox}
\section[\texorpdfstring{\ja{さからう} (sakarau)}{さからう (sakarau)}]{\ja{さからう} (sakarau) — to go against; to disobey}

\label{lesson:JA-C325-sakarau}



\begin{quote}
Before the new one: say the Japanese for an excuse, then the Japanese for to obey.
\end{quote}

\subsection*{You'll want to know: \ja{さからう}}
\textbf{\ja{さからう}} — \emph{sakarau} — \textquotedblleft{}to go against; to disobey\textquotedblright{}.

\begin{quote}
\textbf{\ja{ははに} \ja{さからったので、あとで} \ja{あやまりました。}} — \emph{haha ni sakaratta node, ato de ayamarimashita} — \textquotedblleft{}I went against my mother, so I apologized later.\textquotedblright{}
\end{quote}

Polite \textbf{\ja{さからいます}}, \emph{sakaraimasu}; a godan verb. Past (polite): \textbf{\ja{むすこが} \ja{さからいました。}}, \emph{musuko ga sakaraimashita}.

\begin{cousinweb}
\textbf{\ja{さからう}} (\emph{sakarau}) is a native Japanese word (wago), inherited rather than borrowed from Chinese, and its older origin is not settled. Its kanji was attached to it later as a kun'yomi, a native reading given to a Chinese character (JMdict, KANJIDIC2).
\end{cousinweb}

\begin{grammarlens}[title={What you've built}]
One more word for giving reasons and opinions.
\end{grammarlens}

\subsection*{Your turn}
\begin{itemize}
\raggedright
  \item \emph{Say it:} \emph{sakarau}
  \item \emph{Say it:} \emph{sakarau}, once more
  \item \emph{Say it:} \emph{haha ni sakaratta node, ato de ayamarimashita} — I went against my mother, so I apologized later.
  \item \emph{Recall:} say the Japanese for an excuse, then the Japanese for to obey, then say \emph{sakarau} again
\end{itemize}

\begin{tcolorbox}[breakable,colback=teal!4,colframe=teal!35!black,title={Before you move on}]
What does \ja{さからう} mean? (\textquotedblleft{}To go against; to disobey\textquotedblright{}.) Say it once more.
\end{tcolorbox}
\section[\texorpdfstring{\ja{うけとめる} (uketomeru)}{うけとめる (uketomeru)}]{\ja{うけとめる} (uketomeru) — to take to heart; to catch}

\label{lesson:JA-C325-uketomeru}



\begin{quote}
Before the new one: say the Japanese for to obey, then the Japanese for to go against.
\end{quote}

\subsection*{You'll want to know: \ja{うけとめる}}
\textbf{\ja{うけとめる}} — \emph{uketomeru} — \textquotedblleft{}to take to heart; to catch\textquotedblright{}.

\begin{quote}
\textbf{\ja{いけんが} \ja{ただしいと} \ja{おもったので、ちちは} \ja{うけとめました。}} — \emph{iken ga tadashii to omotta node, chichi wa uketomemashita} — \textquotedblleft{}My father thought the point was fair, so he took it to heart.\textquotedblright{}
\end{quote}

Polite \textbf{\ja{うけとめます}}, \emph{uketomemasu}; an ichidan verb. Past (polite): \textbf{\ja{ちちが} \ja{うけとめました。}}, \emph{chichi ga uketomemashita}.

\begin{cousinweb}
\textbf{\ja{うけとめる}} (\emph{uketomeru}) is a native Japanese word (wago), inherited rather than borrowed from Chinese: it is built from \emph{ukeru}, \textquotedblleft{}to receive\textquotedblright{} and \emph{tomeru}, \textquotedblleft{}to stop\textquotedblright{}, and its kanji were attached later as kun'yomi (JMdict).
\end{cousinweb}

\begin{grammarlens}[title={What you've built}]
One more word for giving reasons and opinions.
\end{grammarlens}

\subsection*{Your turn}
\begin{itemize}
\raggedright
  \item \emph{Say it:} \emph{uketomeru}
  \item \emph{Say it:} \emph{uketomeru}, once more
  \item \emph{Say it:} \emph{iken ga tadashii to omotta node, chichi wa uketomemashita} — My father thought the point was fair, so he took it to heart.
  \item \emph{Recall:} say the Japanese for to obey, then the Japanese for to go against, then say \emph{uketomeru} again
\end{itemize}

\begin{tcolorbox}[breakable,colback=teal!4,colframe=teal!35!black,title={Before you move on}]
What does \ja{うけとめる} mean? (\textquotedblleft{}To take to heart; to catch\textquotedblright{}.) Say it once more.
\end{tcolorbox}
\section[\texorpdfstring{\ja{けっきょく} (kekkyoku)}{けっきょく (kekkyoku)}]{\ja{けっきょく} (kekkyoku) — in the end; after all}

\label{lesson:JA-C325-kekkyoku}



\begin{quote}
Before the new one: say the Japanese for to go against, then the Japanese for to take to heart.
\end{quote}

\subsection*{You'll want to know: \ja{けっきょく}}
\textbf{\ja{けっきょく}} — \emph{kekkyoku} — \textquotedblleft{}in the end; after all\textquotedblright{}.

\begin{quote}
\textbf{\ja{あめが} \ja{ふったので、けっきょく} \ja{いきませんでした。}} — \emph{ame ga futta node, kekkyoku ikimasen deshita} — \textquotedblleft{}It rained, so in the end we did not go.\textquotedblright{}
\end{quote}

\begin{cousinweb}
\textbf{\ja{けっきょく}} (\emph{kekkyoku}) is Sino-Japanese (kango): two Chinese characters read with their on'yomi, readings borrowed from Chinese. They are \emph{ketsu}, \textquotedblleft{}tie\textquotedblright{} (Middle Chinese \emph{ket}) and \emph{kyoku}, \textquotedblleft{}conclusion\textquotedblright{} (\emph{gjowk}) (KANJIDIC2; Middle Chinese in Baxter's transcription). The two join as \emph{kek-kyoku}.
\end{cousinweb}

\begin{grammarlens}[title={What you've built}]
One more word for giving reasons and opinions.
\end{grammarlens}

\subsection*{Your turn}
\begin{itemize}
\raggedright
  \item \emph{Say it:} \emph{kekkyoku}
  \item \emph{Say it:} \emph{kekkyoku}, once more
  \item \emph{Say it:} \emph{ame ga futta node, kekkyoku ikimasen deshita} — It rained, so in the end we did not go.
  \item \emph{Recall:} say the Japanese for to go against, then the Japanese for to take to heart, then say \emph{kekkyoku} again
\end{itemize}

\begin{tcolorbox}[breakable,colback=teal!4,colframe=teal!35!black,title={Before you move on}]
What does \ja{けっきょく} mean? (\textquotedblleft{}In the end; after all\textquotedblright{}.) Say it once more.
\end{tcolorbox}
\section[\texorpdfstring{\ja{じゆう} (jiyū)}{じゆう (jiyū)}]{\ja{じゆう} (jiyū) — freedom; free}

\label{lesson:JA-C325-jiyuu}



\begin{quote}
Before the new one: say the Japanese for to take to heart, then the Japanese for after all.
\end{quote}

\subsection*{You'll want to know: \ja{じゆう}}
\textbf{\ja{じゆう}} — \emph{jiyū} — \textquotedblleft{}freedom; free\textquotedblright{}.

\begin{quote}
\textbf{\ja{この} \ja{がっこうは} \ja{ふくが} \ja{じゆうです。}} — \emph{kono gakkō wa fuku ga jiyū desu} — \textquotedblleft{}At this school, what you wear is up to you.\textquotedblright{}
\end{quote}

A na-adjective: \ja{な} before a noun.

\begin{cousinweb}
\textbf{\ja{じゆう}} (\emph{jiyū}) is Sino-Japanese (kango): two Chinese characters read with their on'yomi, readings borrowed from Chinese. They are \emph{ji}, \textquotedblleft{}oneself\textquotedblright{} (Middle Chinese \emph{dzijH}) and \emph{yū}, \textquotedblleft{}a reason\textquotedblright{} (\emph{yuw}) (KANJIDIC2; Middle Chinese in Baxter's transcription).
\end{cousinweb}

\begin{grammarlens}[title={What you've built}]
One more word for giving reasons and opinions.
\end{grammarlens}

\subsection*{Your turn}
\begin{itemize}
\raggedright
  \item \emph{Say it:} \emph{jiyū}
  \item \emph{Say it:} \emph{jiyū}, once more
  \item \emph{Say it:} \emph{kono gakkō wa fuku ga jiyū desu} — At this school, what you wear is up to you.
  \item \emph{Recall:} say the Japanese for to take to heart, then the Japanese for after all, then say \emph{jiyū} again
\end{itemize}

\begin{tcolorbox}[breakable,colback=teal!4,colframe=teal!35!black,title={Before you move on}]
What does \ja{じゆう} mean? (\textquotedblleft{}Freedom; free\textquotedblright{}.) Say it once more.
\end{tcolorbox}
